\section{Conclusion}

We began by asking whether foundation models changed how we measure AI progress. Our answer is quantitative and definitive: yes, and we can date it precisely.

Using PELT change-point detection on seven years of data, we identified two structural breaks (April 2019, March 2021) with BIC improvement of 17.13. The mechanism is attention reallocation: emergent benchmarks captured +19 percentage points of attention while traditional benchmarks persisted with reduced dominance.

Critically, we discovered that modality dynamics were never unified (pre-2020 $r = -0.131$), refuting our divergence hypothesis but yielding the novel finding that the ecosystem was always siloed by modality.

\subsection{Future Directions}

\begin{itemize}
    \item Test whether multimodal foundation models (CLIP, Flamingo) \emph{synchronized} modality dynamics
    \item Extend analysis pre-2018 for earlier transitions
    \item Implement portfolio churn analysis (benchmark ranking volatility)
    \item Compare PELT results against alternative change-point methods for robustness
\end{itemize}

The rise of foundation models transformed not just what AI can do, but how we measure what AI can do.
